% s1_intro.tex — Introduction
% Paper: "Tokenized Housing and Lifecycle Portfolio Choice:
%         A Decoupling of Location from Housing Exposure"
%
% Source: paper/outline_v4.md §1 (sections 1.1–1.4)
%         docs/welfare_decomp_v4.md §5–6 (contribution framing)
%         docs/calibration_v3.md (empirical anchors)
%
% Status: skeleton complete; [PLACEHOLDER] marks require server1 numerical results.
%         Fill in before submission: CEV headline, channel magnitudes, H1 evidence.
%
% Compilation: % s1_intro.tex — Introduction
% Paper: "Tokenized Housing and Lifecycle Portfolio Choice:
%         A Decoupling of Location from Housing Exposure"
%
% Source: paper/outline_v4.md §1 (sections 1.1–1.4)
%         docs/welfare_decomp_v4.md §5–6 (contribution framing)
%         docs/calibration_v3.md (empirical anchors)
%
% Status: skeleton complete; [PLACEHOLDER] marks require server1 numerical results.
%         Fill in before submission: CEV headline, channel magnitudes, H1 evidence.
%
% Compilation: % s1_intro.tex — Introduction
% Paper: "Tokenized Housing and Lifecycle Portfolio Choice:
%         A Decoupling of Location from Housing Exposure"
%
% Source: paper/outline_v4.md §1 (sections 1.1–1.4)
%         docs/welfare_decomp_v4.md §5–6 (contribution framing)
%         docs/calibration_v3.md (empirical anchors)
%
% Status: skeleton complete; [PLACEHOLDER] marks require server1 numerical results.
%         Fill in before submission: CEV headline, channel magnitudes, H1 evidence.
%
% Compilation: % s1_intro.tex — Introduction
% Paper: "Tokenized Housing and Lifecycle Portfolio Choice:
%         A Decoupling of Location from Housing Exposure"
%
% Source: paper/outline_v4.md §1 (sections 1.1–1.4)
%         docs/welfare_decomp_v4.md §5–6 (contribution framing)
%         docs/calibration_v3.md (empirical anchors)
%
% Status: skeleton complete; [PLACEHOLDER] marks require server1 numerical results.
%         Fill in before submission: CEV headline, channel magnitudes, H1 evidence.
%
% Compilation: \input{sections/s1_intro} from main.tex
%              Requires amsmath, hyperref, natbib (or biblatex).

% ─────────────────────────────────────────────────────────────────────────────
\section{Introduction}
\label{sec:intro}
% ─────────────────────────────────────────────────────────────────────────────

American households relocate frequently.
The Panel Study of Income Dynamics (PSID) records inter-MSA mobility rates of
5--7\% per year among working-age adults, implying that more than 70\% of
households move at least once between ages~25 and~65.
Each relocation under traditional homeownership forces a \emph{sell-and-buy
round-trip}: the household sells the current residence, incurring a broker
commission and closing costs of roughly 6\% of home value, then purchases at
the new location, adding origination fees, title insurance, and inspection
costs of a further 2--3\%.
The aggregate round-trip transaction cost is approximately 8--10\% of home
value per relocation event \citep{NAR2023, CFPB2022}.

These frictions impose a compounding welfare loss over the lifecycle.
A household that relocates four times between ages~25 and~65~---~a plausible
modal count under PSID mobility rates---~faces cumulative round-trip costs
of 32--40\% of a home's value, equivalent to multiple years of mortgage
payments forgone.
Beyond the direct cost, each forced sale severs the household's
\emph{location-specific housing-market exposure}: upon arriving at
location~$B$, the former location-$A$ owner no longer participates in
location~$A$'s price appreciation, even if labor, family, or human-capital
ties make future return migration likely.
Diversified real estate investment trusts (REITs) aggregate at the
portfolio level and provide no location-specific hedge.

Fractional residential housing tokens---blockchain-registered claims to
location-specific property cash flows, exemplified by RealT-class
platforms---offer a structural remedy.
A tokenized holding at location~$A$ can be \emph{retained across
relocation} to location~$B$ at a transfer cost of approximately 1\%,
versus the 8.5\% round-trip for direct ownership.
More importantly, a forward-looking household at location~$A$ who
anticipates a possible future relocation to~$B$ can \emph{pre-accumulate}
fractional location-$B$ tokens incrementally, paying the 2.5\% buying cost
in small installments rather than as a lump-sum at the moment of arrival.
This pre-accumulation motive is the structurally novel mechanism that
distinguishes token portability from both direct ownership and diversified REITs:
the household \emph{decouples its geographic location from its housing-market exposure}.

This paper provides the first lifecycle welfare quantification of
location-decoupling.
We build a two-location lifecycle portfolio model with stochastic relocation
shocks, location-correlated house-price processes, and three ownership regimes:
rent-only (E0), traditional binary homeownership that forces a full sell-and-buy
at every relocation (E1\_2L), and continuous fractional token ownership that
is portable across locations and priced correctly via a 6D state space that
tracks prior-period token holdings (E2\_2L).
The 6D state---which includes last period's token positions $(x_{A,\text{prev}},
x_{B,\text{prev}})$ as explicit state variables---enables the model to price
the per-period buying cost correctly and to generate the forward-looking
pre-accumulation motive that is the paper's central mechanism.

Our headline result is:
\begin{equation}
    \operatorname{CEV}(E2\_2L \text{ vs } E1\_2L) = \textsc{[x.xx\%]},
    \label{eq:headline_cev}
\end{equation}
at baseline calibration ($p_{\text{reloc}} = 0.06$, $\tau_{\text{buy}} = 2.5\%$,
$\rho_{AB} = 0.50$).
[\textsc{Placeholder: fill in once server1 runs confirm.}]
This gain decomposes into three channels:
\begin{enumerate}
    \item \emph{Avoided-transaction-cost channel}: \textsc{[x.x\%]}.
        At each relocation event, the E1\_2L household pays
        $\tau_{\text{sell}} + \tau_{\text{buy}} \approx 8.5\%$;
        the E2\_2L household retains tokens at $\tau_{\text{token}} \approx 1\%$.
        This channel is unique to token portability and has no analog in
        \citet{Liu2021}, who abstracts from mobility.

    \item \emph{Maintained-hedge channel}: \textsc{[x.x\%]}.
        A household at location~$B$ who holds $x_{A,\text{prev}} > 0$ continues
        to participate in location~$A$'s price appreciation.
        The mechanism activates through the pre-accumulation motive enabled by
        the 6D state: the household optimally builds $x_B > 0$ at location~$A$
        to reduce future buying costs on relocation (the ``pre-buy hedge'').
        [\textsc{Placeholder: confirm sign via $\bar{x}_B > 0$ at $\ell = A$
        in server1 output (H1 test).}]

    \item \emph{Continuous-ownership channel}: \textsc{[x.x\%]}.
        Continuous fractional ownership of the occupied residence provides
        rent saving of $\delta_{\text{own}} = \rho - m = 4\%$ per unit held.
        This channel is common to \citet{Liu2021} and replicates his
        MHS-relaxation result within our model.
\end{enumerate}
Channels~(1) and~(2) are the paper's structural contribution beyond
\citet{Liu2021}: token portability uniquely provides (1) because direct
owners face forced sales, and REITs provide no location-specific claims
for channel~(2).
The cross-term between channels is empirically negligible (\textsc{[x.xx\%]})
as in our v3 channel-decomposition benchmark, confirming additive separability.

\paragraph{Falsification structure.}
We pre-register three mechanism tests that must pass for the token-portability
claim to be credible.
(r)~Setting $p_{\text{reloc}} = 0$ shuts off the mobility channel;
    channels~(1) and~(2) must collapse to zero.
(m)~Setting $\rho_{AB} \to 1$ eliminates diversification benefit;
    the maintained-hedge channel must attenuate monotonically.
(q)~Setting $\tau_{\text{buy}} = 0$ eliminates the pre-accumulation motive;
    the pre-buy hedge channel must vanish.
[\textsc{Placeholder: report $p$-values or CEV changes from server1 sweep results.}]

\paragraph{Related literature.}
Our paper connects four strands.

\medskip\noindent
\textbf{Lifecycle housing-portfolio choice.}
\citet{YaoZhang2005} and \citet{Cocco2005} establish the canonical
binary-tenure lifecycle model; we extend by adding a second location,
stochastic relocation, and token portability.
\citet{KraftMunk2011} and \citet{KMW2018} introduce habit and
adjustment-cost extensions; we add geographic mobility.
\citet{Liu2021} introduces the Mortgage-with-Housing-Supply (MHS)
relaxation: continuous ownership fraction within a single location.
Our continuous-$x$ channel (E2\_2L at a single location) replicates Liu's
mechanism; our paper's novel contribution lies in channels~(1) and~(2)
above, which are outside Liu's single-location framework.

\medskip\noindent
\textbf{Housing as hedge against mobility risk.}
\citet{SinaiSouleles2005} show that homeownership hedges against rent-price
risk for households who expect to remain in a location.
\citet{Davidoff2006} documents income-housing correlation as a source of
labor-market risk for owner-occupiers.
\citet{BFN2014} (Bagliano, Fugazza, and Nicodano) show that housing and
equity exposure interact through labor income in a lifecycle model.
Our contribution is distinct: we model the \emph{loss} of location-specific
hedge at relocation and quantify the welfare value of maintaining it via tokens.

\medskip\noindent
\textbf{Transaction costs in housing markets.}
\citet{FlavinYamashita2002} document that transaction costs constrain
portfolio rebalancing in a housing lifecycle context.
\citet{Han2013} quantifies the lock-in effect of transaction costs on
tenure choice.
\citet{PiazzesiSchneider2016} embed transaction costs in an equilibrium
housing-search model.
Our model treats transaction costs as the \emph{source} of the E1\_2L
welfare loss; token portability is valued precisely because it reduces these costs.

\medskip\noindent
\textbf{Tokenization and digital assets.}
\citet{CongLiWang2021} provide the foundational theory of tokenomics in a
general equilibrium setting.
\citet{Swinkels2023} surveys the institutional structure of residential
real estate tokenization platforms.
Our contribution is a lifecycle welfare quantification \emph{within} a
standard Cocco-class model, grounded in the PSID-/NAR-calibrated parameter
space, rather than a platform mechanism design.

\paragraph{Partial equilibrium caveat.}
All results are partial-equilibrium: house prices are exogenous.
We interpret $\operatorname{CEV}(E2\_2L \text{ vs } E1\_2L)$ as the
household-level welfare value of \emph{access} to tokenized instruments,
not as an aggregate welfare estimate.
A general-equilibrium extension with endogenous prices and adoption externalities
is left for future work.

\paragraph{Road map.}
Section~\ref{sec:model} presents the two-location lifecycle model.
Section~\ref{sec:calibration} describes the calibration.
Section~\ref{sec:results} reports baseline welfare results and the channel
decomposition.
Section~\ref{sec:discussion} discusses the contribution relative to
\citet{Liu2021} and the robustness of the falsification structure.
Section~\ref{sec:conclusion} concludes.
All numerical details are in the appendices.
 from main.tex
%              Requires amsmath, hyperref, natbib (or biblatex).

% ─────────────────────────────────────────────────────────────────────────────
\section{Introduction}
\label{sec:intro}
% ─────────────────────────────────────────────────────────────────────────────

American households relocate frequently.
The Panel Study of Income Dynamics (PSID) records inter-MSA mobility rates of
5--7\% per year among working-age adults, implying that more than 70\% of
households move at least once between ages~25 and~65.
Each relocation under traditional homeownership forces a \emph{sell-and-buy
round-trip}: the household sells the current residence, incurring a broker
commission and closing costs of roughly 6\% of home value, then purchases at
the new location, adding origination fees, title insurance, and inspection
costs of a further 2--3\%.
The aggregate round-trip transaction cost is approximately 8--10\% of home
value per relocation event \citep{NAR2023, CFPB2022}.

These frictions impose a compounding welfare loss over the lifecycle.
A household that relocates four times between ages~25 and~65~---~a plausible
modal count under PSID mobility rates---~faces cumulative round-trip costs
of 32--40\% of a home's value, equivalent to multiple years of mortgage
payments forgone.
Beyond the direct cost, each forced sale severs the household's
\emph{location-specific housing-market exposure}: upon arriving at
location~$B$, the former location-$A$ owner no longer participates in
location~$A$'s price appreciation, even if labor, family, or human-capital
ties make future return migration likely.
Diversified real estate investment trusts (REITs) aggregate at the
portfolio level and provide no location-specific hedge.

Fractional residential housing tokens---blockchain-registered claims to
location-specific property cash flows, exemplified by RealT-class
platforms---offer a structural remedy.
A tokenized holding at location~$A$ can be \emph{retained across
relocation} to location~$B$ at a transfer cost of approximately 1\%,
versus the 8.5\% round-trip for direct ownership.
More importantly, a forward-looking household at location~$A$ who
anticipates a possible future relocation to~$B$ can \emph{pre-accumulate}
fractional location-$B$ tokens incrementally, paying the 2.5\% buying cost
in small installments rather than as a lump-sum at the moment of arrival.
This pre-accumulation motive is the structurally novel mechanism that
distinguishes token portability from both direct ownership and diversified REITs:
the household \emph{decouples its geographic location from its housing-market exposure}.

This paper provides the first lifecycle welfare quantification of
location-decoupling.
We build a two-location lifecycle portfolio model with stochastic relocation
shocks, location-correlated house-price processes, and three ownership regimes:
rent-only (E0), traditional binary homeownership that forces a full sell-and-buy
at every relocation (E1\_2L), and continuous fractional token ownership that
is portable across locations and priced correctly via a 6D state space that
tracks prior-period token holdings (E2\_2L).
The 6D state---which includes last period's token positions $(x_{A,\text{prev}},
x_{B,\text{prev}})$ as explicit state variables---enables the model to price
the per-period buying cost correctly and to generate the forward-looking
pre-accumulation motive that is the paper's central mechanism.

Our headline result is:
\begin{equation}
    \operatorname{CEV}(E2\_2L \text{ vs } E1\_2L) = \textsc{[x.xx\%]},
    \label{eq:headline_cev}
\end{equation}
at baseline calibration ($p_{\text{reloc}} = 0.06$, $\tau_{\text{buy}} = 2.5\%$,
$\rho_{AB} = 0.50$).
[\textsc{Placeholder: fill in once server1 runs confirm.}]
This gain decomposes into three channels:
\begin{enumerate}
    \item \emph{Avoided-transaction-cost channel}: \textsc{[x.x\%]}.
        At each relocation event, the E1\_2L household pays
        $\tau_{\text{sell}} + \tau_{\text{buy}} \approx 8.5\%$;
        the E2\_2L household retains tokens at $\tau_{\text{token}} \approx 1\%$.
        This channel is unique to token portability and has no analog in
        \citet{Liu2021}, who abstracts from mobility.

    \item \emph{Maintained-hedge channel}: \textsc{[x.x\%]}.
        A household at location~$B$ who holds $x_{A,\text{prev}} > 0$ continues
        to participate in location~$A$'s price appreciation.
        The mechanism activates through the pre-accumulation motive enabled by
        the 6D state: the household optimally builds $x_B > 0$ at location~$A$
        to reduce future buying costs on relocation (the ``pre-buy hedge'').
        [\textsc{Placeholder: confirm sign via $\bar{x}_B > 0$ at $\ell = A$
        in server1 output (H1 test).}]

    \item \emph{Continuous-ownership channel}: \textsc{[x.x\%]}.
        Continuous fractional ownership of the occupied residence provides
        rent saving of $\delta_{\text{own}} = \rho - m = 4\%$ per unit held.
        This channel is common to \citet{Liu2021} and replicates his
        MHS-relaxation result within our model.
\end{enumerate}
Channels~(1) and~(2) are the paper's structural contribution beyond
\citet{Liu2021}: token portability uniquely provides (1) because direct
owners face forced sales, and REITs provide no location-specific claims
for channel~(2).
The cross-term between channels is empirically negligible (\textsc{[x.xx\%]})
as in our v3 channel-decomposition benchmark, confirming additive separability.

\paragraph{Falsification structure.}
We pre-register three mechanism tests that must pass for the token-portability
claim to be credible.
(r)~Setting $p_{\text{reloc}} = 0$ shuts off the mobility channel;
    channels~(1) and~(2) must collapse to zero.
(m)~Setting $\rho_{AB} \to 1$ eliminates diversification benefit;
    the maintained-hedge channel must attenuate monotonically.
(q)~Setting $\tau_{\text{buy}} = 0$ eliminates the pre-accumulation motive;
    the pre-buy hedge channel must vanish.
[\textsc{Placeholder: report $p$-values or CEV changes from server1 sweep results.}]

\paragraph{Related literature.}
Our paper connects four strands.

\medskip\noindent
\textbf{Lifecycle housing-portfolio choice.}
\citet{YaoZhang2005} and \citet{Cocco2005} establish the canonical
binary-tenure lifecycle model; we extend by adding a second location,
stochastic relocation, and token portability.
\citet{KraftMunk2011} and \citet{KMW2018} introduce habit and
adjustment-cost extensions; we add geographic mobility.
\citet{Liu2021} introduces the Mortgage-with-Housing-Supply (MHS)
relaxation: continuous ownership fraction within a single location.
Our continuous-$x$ channel (E2\_2L at a single location) replicates Liu's
mechanism; our paper's novel contribution lies in channels~(1) and~(2)
above, which are outside Liu's single-location framework.

\medskip\noindent
\textbf{Housing as hedge against mobility risk.}
\citet{SinaiSouleles2005} show that homeownership hedges against rent-price
risk for households who expect to remain in a location.
\citet{Davidoff2006} documents income-housing correlation as a source of
labor-market risk for owner-occupiers.
\citet{BFN2014} (Bagliano, Fugazza, and Nicodano) show that housing and
equity exposure interact through labor income in a lifecycle model.
Our contribution is distinct: we model the \emph{loss} of location-specific
hedge at relocation and quantify the welfare value of maintaining it via tokens.

\medskip\noindent
\textbf{Transaction costs in housing markets.}
\citet{FlavinYamashita2002} document that transaction costs constrain
portfolio rebalancing in a housing lifecycle context.
\citet{Han2013} quantifies the lock-in effect of transaction costs on
tenure choice.
\citet{PiazzesiSchneider2016} embed transaction costs in an equilibrium
housing-search model.
Our model treats transaction costs as the \emph{source} of the E1\_2L
welfare loss; token portability is valued precisely because it reduces these costs.

\medskip\noindent
\textbf{Tokenization and digital assets.}
\citet{CongLiWang2021} provide the foundational theory of tokenomics in a
general equilibrium setting.
\citet{Swinkels2023} surveys the institutional structure of residential
real estate tokenization platforms.
Our contribution is a lifecycle welfare quantification \emph{within} a
standard Cocco-class model, grounded in the PSID-/NAR-calibrated parameter
space, rather than a platform mechanism design.

\paragraph{Partial equilibrium caveat.}
All results are partial-equilibrium: house prices are exogenous.
We interpret $\operatorname{CEV}(E2\_2L \text{ vs } E1\_2L)$ as the
household-level welfare value of \emph{access} to tokenized instruments,
not as an aggregate welfare estimate.
A general-equilibrium extension with endogenous prices and adoption externalities
is left for future work.

\paragraph{Road map.}
Section~\ref{sec:model} presents the two-location lifecycle model.
Section~\ref{sec:calibration} describes the calibration.
Section~\ref{sec:results} reports baseline welfare results and the channel
decomposition.
Section~\ref{sec:discussion} discusses the contribution relative to
\citet{Liu2021} and the robustness of the falsification structure.
Section~\ref{sec:conclusion} concludes.
All numerical details are in the appendices.
 from main.tex
%              Requires amsmath, hyperref, natbib (or biblatex).

% ─────────────────────────────────────────────────────────────────────────────
\section{Introduction}
\label{sec:intro}
% ─────────────────────────────────────────────────────────────────────────────

American households relocate frequently.
The Panel Study of Income Dynamics (PSID) records inter-MSA mobility rates of
5--7\% per year among working-age adults, implying that more than 70\% of
households move at least once between ages~25 and~65.
Each relocation under traditional homeownership forces a \emph{sell-and-buy
round-trip}: the household sells the current residence, incurring a broker
commission and closing costs of roughly 6\% of home value, then purchases at
the new location, adding origination fees, title insurance, and inspection
costs of a further 2--3\%.
The aggregate round-trip transaction cost is approximately 8--10\% of home
value per relocation event \citep{NAR2023, CFPB2022}.

These frictions impose a compounding welfare loss over the lifecycle.
A household that relocates four times between ages~25 and~65~---~a plausible
modal count under PSID mobility rates---~faces cumulative round-trip costs
of 32--40\% of a home's value, equivalent to multiple years of mortgage
payments forgone.
Beyond the direct cost, each forced sale severs the household's
\emph{location-specific housing-market exposure}: upon arriving at
location~$B$, the former location-$A$ owner no longer participates in
location~$A$'s price appreciation, even if labor, family, or human-capital
ties make future return migration likely.
Diversified real estate investment trusts (REITs) aggregate at the
portfolio level and provide no location-specific hedge.

Fractional residential housing tokens---blockchain-registered claims to
location-specific property cash flows, exemplified by RealT-class
platforms---offer a structural remedy.
A tokenized holding at location~$A$ can be \emph{retained across
relocation} to location~$B$ at a transfer cost of approximately 1\%,
versus the 8.5\% round-trip for direct ownership.
More importantly, a forward-looking household at location~$A$ who
anticipates a possible future relocation to~$B$ can \emph{pre-accumulate}
fractional location-$B$ tokens incrementally, paying the 2.5\% buying cost
in small installments rather than as a lump-sum at the moment of arrival.
This pre-accumulation motive is the structurally novel mechanism that
distinguishes token portability from both direct ownership and diversified REITs:
the household \emph{decouples its geographic location from its housing-market exposure}.

This paper provides the first lifecycle welfare quantification of
location-decoupling.
We build a two-location lifecycle portfolio model with stochastic relocation
shocks, location-correlated house-price processes, and three ownership regimes:
rent-only (E0), traditional binary homeownership that forces a full sell-and-buy
at every relocation (E1\_2L), and continuous fractional token ownership that
is portable across locations and priced correctly via a 6D state space that
tracks prior-period token holdings (E2\_2L).
The 6D state---which includes last period's token positions $(x_{A,\text{prev}},
x_{B,\text{prev}})$ as explicit state variables---enables the model to price
the per-period buying cost correctly and to generate the forward-looking
pre-accumulation motive that is the paper's central mechanism.

Our headline result is:
\begin{equation}
    \operatorname{CEV}(E2\_2L \text{ vs } E1\_2L) = \textsc{[x.xx\%]},
    \label{eq:headline_cev}
\end{equation}
at baseline calibration ($p_{\text{reloc}} = 0.06$, $\tau_{\text{buy}} = 2.5\%$,
$\rho_{AB} = 0.50$).
[\textsc{Placeholder: fill in once server1 runs confirm.}]
This gain decomposes into three channels:
\begin{enumerate}
    \item \emph{Avoided-transaction-cost channel}: \textsc{[x.x\%]}.
        At each relocation event, the E1\_2L household pays
        $\tau_{\text{sell}} + \tau_{\text{buy}} \approx 8.5\%$;
        the E2\_2L household retains tokens at $\tau_{\text{token}} \approx 1\%$.
        This channel is unique to token portability and has no analog in
        \citet{Liu2021}, who abstracts from mobility.

    \item \emph{Maintained-hedge channel}: \textsc{[x.x\%]}.
        A household at location~$B$ who holds $x_{A,\text{prev}} > 0$ continues
        to participate in location~$A$'s price appreciation.
        The mechanism activates through the pre-accumulation motive enabled by
        the 6D state: the household optimally builds $x_B > 0$ at location~$A$
        to reduce future buying costs on relocation (the ``pre-buy hedge'').
        [\textsc{Placeholder: confirm sign via $\bar{x}_B > 0$ at $\ell = A$
        in server1 output (H1 test).}]

    \item \emph{Continuous-ownership channel}: \textsc{[x.x\%]}.
        Continuous fractional ownership of the occupied residence provides
        rent saving of $\delta_{\text{own}} = \rho - m = 4\%$ per unit held.
        This channel is common to \citet{Liu2021} and replicates his
        MHS-relaxation result within our model.
\end{enumerate}
Channels~(1) and~(2) are the paper's structural contribution beyond
\citet{Liu2021}: token portability uniquely provides (1) because direct
owners face forced sales, and REITs provide no location-specific claims
for channel~(2).
The cross-term between channels is empirically negligible (\textsc{[x.xx\%]})
as in our v3 channel-decomposition benchmark, confirming additive separability.

\paragraph{Falsification structure.}
We pre-register three mechanism tests that must pass for the token-portability
claim to be credible.
(r)~Setting $p_{\text{reloc}} = 0$ shuts off the mobility channel;
    channels~(1) and~(2) must collapse to zero.
(m)~Setting $\rho_{AB} \to 1$ eliminates diversification benefit;
    the maintained-hedge channel must attenuate monotonically.
(q)~Setting $\tau_{\text{buy}} = 0$ eliminates the pre-accumulation motive;
    the pre-buy hedge channel must vanish.
[\textsc{Placeholder: report $p$-values or CEV changes from server1 sweep results.}]

\paragraph{Related literature.}
Our paper connects four strands.

\medskip\noindent
\textbf{Lifecycle housing-portfolio choice.}
\citet{YaoZhang2005} and \citet{Cocco2005} establish the canonical
binary-tenure lifecycle model; we extend by adding a second location,
stochastic relocation, and token portability.
\citet{KraftMunk2011} and \citet{KMW2018} introduce habit and
adjustment-cost extensions; we add geographic mobility.
\citet{Liu2021} introduces the Mortgage-with-Housing-Supply (MHS)
relaxation: continuous ownership fraction within a single location.
Our continuous-$x$ channel (E2\_2L at a single location) replicates Liu's
mechanism; our paper's novel contribution lies in channels~(1) and~(2)
above, which are outside Liu's single-location framework.

\medskip\noindent
\textbf{Housing as hedge against mobility risk.}
\citet{SinaiSouleles2005} show that homeownership hedges against rent-price
risk for households who expect to remain in a location.
\citet{Davidoff2006} documents income-housing correlation as a source of
labor-market risk for owner-occupiers.
\citet{BFN2014} (Bagliano, Fugazza, and Nicodano) show that housing and
equity exposure interact through labor income in a lifecycle model.
Our contribution is distinct: we model the \emph{loss} of location-specific
hedge at relocation and quantify the welfare value of maintaining it via tokens.

\medskip\noindent
\textbf{Transaction costs in housing markets.}
\citet{FlavinYamashita2002} document that transaction costs constrain
portfolio rebalancing in a housing lifecycle context.
\citet{Han2013} quantifies the lock-in effect of transaction costs on
tenure choice.
\citet{PiazzesiSchneider2016} embed transaction costs in an equilibrium
housing-search model.
Our model treats transaction costs as the \emph{source} of the E1\_2L
welfare loss; token portability is valued precisely because it reduces these costs.

\medskip\noindent
\textbf{Tokenization and digital assets.}
\citet{CongLiWang2021} provide the foundational theory of tokenomics in a
general equilibrium setting.
\citet{Swinkels2023} surveys the institutional structure of residential
real estate tokenization platforms.
Our contribution is a lifecycle welfare quantification \emph{within} a
standard Cocco-class model, grounded in the PSID-/NAR-calibrated parameter
space, rather than a platform mechanism design.

\paragraph{Partial equilibrium caveat.}
All results are partial-equilibrium: house prices are exogenous.
We interpret $\operatorname{CEV}(E2\_2L \text{ vs } E1\_2L)$ as the
household-level welfare value of \emph{access} to tokenized instruments,
not as an aggregate welfare estimate.
A general-equilibrium extension with endogenous prices and adoption externalities
is left for future work.

\paragraph{Road map.}
Section~\ref{sec:model} presents the two-location lifecycle model.
Section~\ref{sec:calibration} describes the calibration.
Section~\ref{sec:results} reports baseline welfare results and the channel
decomposition.
Section~\ref{sec:discussion} discusses the contribution relative to
\citet{Liu2021} and the robustness of the falsification structure.
Section~\ref{sec:conclusion} concludes.
All numerical details are in the appendices.
 from main.tex
%              Requires amsmath, hyperref, natbib (or biblatex).

% ─────────────────────────────────────────────────────────────────────────────
\section{Introduction}
\label{sec:intro}
% ─────────────────────────────────────────────────────────────────────────────

American households relocate frequently.
The Panel Study of Income Dynamics (PSID) records inter-MSA mobility rates of
5--7\% per year among working-age adults, implying that more than 70\% of
households move at least once between ages~25 and~65.
Each relocation under traditional homeownership forces a \emph{sell-and-buy
round-trip}: the household sells the current residence, incurring a broker
commission and closing costs of roughly 6\% of home value, then purchases at
the new location, adding origination fees, title insurance, and inspection
costs of a further 2--3\%.
The aggregate round-trip transaction cost is approximately 8--10\% of home
value per relocation event \citep{NAR2023, CFPB2022}.

These frictions impose a compounding welfare loss over the lifecycle.
A household that relocates four times between ages~25 and~65~---~a plausible
modal count under PSID mobility rates---~faces cumulative round-trip costs
of 32--40\% of a home's value, equivalent to multiple years of mortgage
payments forgone.
Beyond the direct cost, each forced sale severs the household's
\emph{location-specific housing-market exposure}: upon arriving at
location~$B$, the former location-$A$ owner no longer participates in
location~$A$'s price appreciation, even if labor, family, or human-capital
ties make future return migration likely.
Diversified real estate investment trusts (REITs) aggregate at the
portfolio level and provide no location-specific hedge.

Fractional residential housing tokens---blockchain-registered claims to
location-specific property cash flows, exemplified by RealT-class
platforms---offer a structural remedy.
A tokenized holding at location~$A$ can be \emph{retained across
relocation} to location~$B$ at a transfer cost of approximately 1\%,
versus the 8.5\% round-trip for direct ownership.
More importantly, a forward-looking household at location~$A$ who
anticipates a possible future relocation to~$B$ can \emph{pre-accumulate}
fractional location-$B$ tokens incrementally, paying the 2.5\% buying cost
in small installments rather than as a lump-sum at the moment of arrival.
This pre-accumulation motive is the structurally novel mechanism that
distinguishes token portability from both direct ownership and diversified REITs:
the household \emph{decouples its geographic location from its housing-market exposure}.

This paper provides the first lifecycle welfare quantification of
location-decoupling.
We build a two-location lifecycle portfolio model with stochastic relocation
shocks, location-correlated house-price processes, and three ownership regimes:
rent-only (E0), traditional binary homeownership that forces a full sell-and-buy
at every relocation (E1\_2L), and continuous fractional token ownership that
is portable across locations and priced correctly via a 6D state space that
tracks prior-period token holdings (E2\_2L).
The 6D state---which includes last period's token positions $(x_{A,\text{prev}},
x_{B,\text{prev}})$ as explicit state variables---enables the model to price
the per-period buying cost correctly and to generate the forward-looking
pre-accumulation motive that is the paper's central mechanism.

Our headline result is:
\begin{equation}
    \operatorname{CEV}(E2\_2L \text{ vs } E1\_2L) = \textsc{[x.xx\%]},
    \label{eq:headline_cev}
\end{equation}
at baseline calibration ($p_{\text{reloc}} = 0.06$, $\tau_{\text{buy}} = 2.5\%$,
$\rho_{AB} = 0.50$).
[\textsc{Placeholder: fill in once server1 runs confirm.}]
This gain decomposes into three channels:
\begin{enumerate}
    \item \emph{Avoided-transaction-cost channel}: \textsc{[x.x\%]}.
        At each relocation event, the E1\_2L household pays
        $\tau_{\text{sell}} + \tau_{\text{buy}} \approx 8.5\%$;
        the E2\_2L household retains tokens at $\tau_{\text{token}} \approx 1\%$.
        This channel is unique to token portability and has no analog in
        \citet{Liu2021}, who abstracts from mobility.

    \item \emph{Maintained-hedge channel}: \textsc{[x.x\%]}.
        A household at location~$B$ who holds $x_{A,\text{prev}} > 0$ continues
        to participate in location~$A$'s price appreciation.
        The mechanism activates through the pre-accumulation motive enabled by
        the 6D state: the household optimally builds $x_B > 0$ at location~$A$
        to reduce future buying costs on relocation (the ``pre-buy hedge'').
        [\textsc{Placeholder: confirm sign via $\bar{x}_B > 0$ at $\ell = A$
        in server1 output (H1 test).}]

    \item \emph{Continuous-ownership channel}: \textsc{[x.x\%]}.
        Continuous fractional ownership of the occupied residence provides
        rent saving of $\delta_{\text{own}} = \rho - m = 4\%$ per unit held.
        This channel is common to \citet{Liu2021} and replicates his
        MHS-relaxation result within our model.
\end{enumerate}
Channels~(1) and~(2) are the paper's structural contribution beyond
\citet{Liu2021}: token portability uniquely provides (1) because direct
owners face forced sales, and REITs provide no location-specific claims
for channel~(2).
The cross-term between channels is empirically negligible (\textsc{[x.xx\%]})
as in our v3 channel-decomposition benchmark, confirming additive separability.

\paragraph{Falsification structure.}
We pre-register three mechanism tests that must pass for the token-portability
claim to be credible.
(r)~Setting $p_{\text{reloc}} = 0$ shuts off the mobility channel;
    channels~(1) and~(2) must collapse to zero.
(m)~Setting $\rho_{AB} \to 1$ eliminates diversification benefit;
    the maintained-hedge channel must attenuate monotonically.
(q)~Setting $\tau_{\text{buy}} = 0$ eliminates the pre-accumulation motive;
    the pre-buy hedge channel must vanish.
[\textsc{Placeholder: report $p$-values or CEV changes from server1 sweep results.}]

\paragraph{Related literature.}
Our paper connects four strands.

\medskip\noindent
\textbf{Lifecycle housing-portfolio choice.}
\citet{YaoZhang2005} and \citet{Cocco2005} establish the canonical
binary-tenure lifecycle model; we extend by adding a second location,
stochastic relocation, and token portability.
\citet{KraftMunk2011} and \citet{KMW2018} introduce habit and
adjustment-cost extensions; we add geographic mobility.
\citet{Liu2021} introduces the Mortgage-with-Housing-Supply (MHS)
relaxation: continuous ownership fraction within a single location.
Our continuous-$x$ channel (E2\_2L at a single location) replicates Liu's
mechanism; our paper's novel contribution lies in channels~(1) and~(2)
above, which are outside Liu's single-location framework.

\medskip\noindent
\textbf{Housing as hedge against mobility risk.}
\citet{SinaiSouleles2005} show that homeownership hedges against rent-price
risk for households who expect to remain in a location.
\citet{Davidoff2006} documents income-housing correlation as a source of
labor-market risk for owner-occupiers.
\citet{BFN2014} (Bagliano, Fugazza, and Nicodano) show that housing and
equity exposure interact through labor income in a lifecycle model.
Our contribution is distinct: we model the \emph{loss} of location-specific
hedge at relocation and quantify the welfare value of maintaining it via tokens.

\medskip\noindent
\textbf{Transaction costs in housing markets.}
\citet{FlavinYamashita2002} document that transaction costs constrain
portfolio rebalancing in a housing lifecycle context.
\citet{Han2013} quantifies the lock-in effect of transaction costs on
tenure choice.
\citet{PiazzesiSchneider2016} embed transaction costs in an equilibrium
housing-search model.
Our model treats transaction costs as the \emph{source} of the E1\_2L
welfare loss; token portability is valued precisely because it reduces these costs.

\medskip\noindent
\textbf{Tokenization and digital assets.}
\citet{CongLiWang2021} provide the foundational theory of tokenomics in a
general equilibrium setting.
\citet{Swinkels2023} surveys the institutional structure of residential
real estate tokenization platforms.
Our contribution is a lifecycle welfare quantification \emph{within} a
standard Cocco-class model, grounded in the PSID-/NAR-calibrated parameter
space, rather than a platform mechanism design.

\paragraph{Partial equilibrium caveat.}
All results are partial-equilibrium: house prices are exogenous.
We interpret $\operatorname{CEV}(E2\_2L \text{ vs } E1\_2L)$ as the
household-level welfare value of \emph{access} to tokenized instruments,
not as an aggregate welfare estimate.
A general-equilibrium extension with endogenous prices and adoption externalities
is left for future work.

\paragraph{Road map.}
Section~\ref{sec:model} presents the two-location lifecycle model.
Section~\ref{sec:calibration} describes the calibration.
Section~\ref{sec:results} reports baseline welfare results and the channel
decomposition.
Section~\ref{sec:discussion} discusses the contribution relative to
\citet{Liu2021} and the robustness of the falsification structure.
Section~\ref{sec:conclusion} concludes.
All numerical details are in the appendices.
